\chapter{Introduction et contexte}
\label{chap:intro}

\section{Contexte et problématique}

\subsection{Contexte pédagogique et scientifique}

Ce projet constitue l'évaluation finale du cours d'Interprétation et Compilation
de la Licence~3 Informatique de l'IED, à l'Université Paris~8. Il occupe une
place particulière dans la progression du semestre, car il ne porte pas sur une
notion isolée mais demande au contraire de rassembler en un seul logiciel
fonctionnel l'ensemble des concepts étudiés tout au long de l'année. Les chapitres
précédents avaient abordé séparément l'assembleur et la pile d'appels, la chaîne
de compilation, les automates finis, les expressions régulières, les grammaires,
puis les optimisations. Le projet final consiste à montrer qu'on est capable de
faire tenir tous ces éléments ensemble, de les articuler les uns aux autres, et
d'obtenir au bout un programme qui exécute réellement du code écrit dans un langage
que l'on a soi-même défini.

Il faut rappeler à quel point la traduction des langages occupe une position
centrale en informatique. Aucun programme écrit dans un langage de haut niveau
n'est directement compréhensible par le processeur, qui ne sait exécuter que des
instructions machine élémentaires. Il existe donc toujours, quelque part, un
logiciel chargé de faire le pont entre ce que le programmeur écrit et ce que la
machine exécute. Deux grandes familles de stratégies remplissent ce rôle. La
compilation traduit le programme une fois pour toutes en code machine, qui sera
ensuite exécuté de façon autonome. L'interprétation, à l'inverse, exécute le
programme directement, en parcourant une représentation interne, sans produire de
binaire séparé. Ces deux familles, malgré leur différence de finalité, partagent
exactement la même première moitié de traitement : analyser le texte source, le
décomposer en unités lexicales, vérifier sa conformité à une grammaire, et le
représenter sous une forme structurée. Elles ne divergent qu'ensuite, sur la
manière d'exploiter cette représentation. C'est précisément cette architecture en
deux temps, un tronc commun suivi de traitements spécialisés, que mon projet met
en évidence, d'autant plus nettement qu'il en propose deux exploitations
différentes du même tronc.

\subsection{Identification du problème}

Le problème posé peut s'énoncer simplement : concevoir un langage de programmation
complet, si modeste soit-il, et écrire le logiciel capable de l'exécuter. Le mot
« complet » mérite une précision. Il ne signifie pas riche en fonctionnalités,
mais suffisant pour exprimer n'importe quel calcul, c'est-à-dire Turing-complet.
En pratique, cela impose au minimum de disposer de variables pour mémoriser des
valeurs, de conditionnelles pour choisir entre plusieurs comportements, de boucles
pour répéter des traitements, et d'un mécanisme de composition, ici les fonctions,
pour structurer et réutiliser le code.

Concevoir un tel langage ne se réduit pas à écrire du code C. C'est un travail
préalable de définition qui porte sur trois aspects indissociables, et c'est
seulement une fois ces trois aspects fixés que l'écriture de l'interpréteur prend
son sens. Le premier aspect est la syntaxe concrète, c'est-à-dire l'apparence
qu'aura un programme valide : quels mots-clés, quels opérateurs, quelle
ponctuation. Le deuxième aspect est la grammaire, c'est-à-dire la définition
précise des suites de symboles acceptées et de la manière dont elles s'organisent
en structures imbriquées. Le troisième aspect, enfin, est la sémantique,
c'est-à-dire la signification de chaque construction, ce que fait réellement le
programme lorsqu'on l'exécute. La syntaxe relève de Flex, la grammaire de Bison,
et la sémantique de l'évaluateur que j'ai écrit.

\subsection{Questions directrices}

Plutôt que d'aborder le projet comme une simple commande à satisfaire, j'ai
préféré l'organiser autour de quelques questions concrètes d'ingénierie, qui ont
guidé mes choix et qui structurent ce rapport. La première question concerne les
outils : qu'apporte réellement le recours à Flex et Bison par rapport à une analyse
écrite à la main, et ce gain justifie-t-il la courbe d'apprentissage de ces
outils ? La deuxième question concerne la représentation interne : comment
représenter un programme en mémoire de manière à pouvoir aussi bien l'évaluer,
l'optimiser, voire le compiler, sans avoir à tout réécrire pour chaque usage ? La
troisième question touche à la sémantique des fonctions : comment gérer
correctement la portée des variables et la récursion, qui sont les points les plus
délicats d'un langage à fonctions ? La dernière question, plus ouverte, est celle
des limites : jusqu'où peut-on pousser un projet de ce type avant de buter sur des
obstacles de fond, qu'il s'agisse de performance, de profondeur de récursion ou de
typage ?

\section{Objectifs et contributions}

\subsection{Objectifs}

L'objectif principal est de livrer un interpréteur \ilang{} qui soit à la fois
fonctionnel, robuste et fidèle au cahier des charges. Cet objectif se décompose en
plusieurs sous-objectifs concrets. Il faut d'abord écrire un analyseur lexical qui
reconnaisse toutes les unités lexicales du langage, des nombres aux mots-clés en
passant par les commentaires. Il faut ensuite écrire une grammaire LALR(1) qui
soit non ambiguë et qui construise au fil de l'analyse l'arbre syntaxique. Il faut
implémenter un évaluateur capable de traiter les variables, les expressions, les
structures de contrôle, les fonctions récursives et les entrées-sorties. Il faut
mettre en place une table des symboles correcte du point de vue de la portée. Il
faut réaliser les deux passes d'optimisation demandées. Il faut traiter proprement
les erreurs courantes en les localisant à la ligne. Enfin, et ce dernier point a
compté autant que les autres à mes yeux, il faut atteindre un certain niveau de
qualité d'ingénierie : une compilation sans le moindre avertissement, une grammaire
sans conflit, et une exécution sans fuite mémoire.

\subsection{Contributions}

Au-delà de ce que demandait strictement la commande, mon travail apporte trois
contributions que je revendique. La première est un ensemble de trois extensions
du langage. J'ai ajouté un type flottant, conçu de façon à cohabiter avec les
entiers sans compromettre la sémantique entière d'origine, ce qui était la vraie
difficulté. J'ai ajouté les instructions de rupture de boucle \code{arrêter} et
\code{continuer}. J'ai enfin systématisé la localisation des erreurs à la ligne.
La deuxième contribution est un second back-end, qui génère du code assembleur
i386 à partir du même arbre syntaxique. Je l'ai développé par curiosité, en sachant
qu'il dépassait le périmètre du projet, mais il démontre concrètement que le
front-end que j'avais écrit était réutilisable tel quel pour produire du code
natif, ce qui valide a posteriori toute l'architecture. La troisième contribution
est une démarche de validation que j'ai voulue sérieuse, comprenant des tests
fonctionnels, des tests d'erreurs, une campagne de cas limites, une analyse de la
mémoire et des mesures de performance comparant l'interprétation au code natif.

\subsection{Cadrage du projet par les cinq questions}

Pour cadrer clairement le travail, il est utile de répondre explicitement aux cinq
questions classiques de la conduite de projet, souvent appelées les cinq W. La
question du \emph{quoi} a déjà été posée : il s'agit d'un interpréteur pour le
mini-langage impératif \ilang, accompagné de ses programmes de démonstration et de
sa documentation. La question du \emph{pourquoi} renvoie à la finalité
pédagogique : il s'agit de démontrer la maîtrise de la chaîne de compilation
complète, ce qui est l'objectif assigné au projet final du chapitre~10. La question
du \emph{qui} est vite réglée, puisqu'il s'agit d'un travail individuel évalué,
réalisé par un seul étudiant. La question du \emph{quand} renvoie à la période
allouée en fin de semestre, dont la répartition est présentée sous forme de
diagramme de Gantt au chapitre~\ref{chap:conception}. La question du \emph{où},
enfin, désigne l'environnement de travail, à savoir un système GNU/Linux exécuté
sous WSL, avec la chaîne d'outils GCC, Flex et Bison.

\subsection{Budget, temps et coûts}

S'agissant d'un projet académique individuel, la question financière se règle en
une phrase : le coût direct est nul. Tous les outils que j'ai employés, du
compilateur GCC aux générateurs Flex et Bison, en passant par Make, Valgrind et la
distribution \TeX{} Live qui a servi à produire ce rapport, sont des logiciels
libres et gratuits. La seule ressource réellement consommée est le temps. J'ai
réparti cet effort de manière progressive, en commençant par la conception de la
grammaire, puis l'écriture du c\oe ur du langage, ensuite la table des symboles et
les optimisations, puis une phase substantielle de débogage et de tests, et enfin
les extensions avant la rédaction. Cette répartition est détaillée et justifiée au
chapitre~\ref{chap:conception}, où elle est présentée sous forme de planning.

\section{Réinvestissement des acquis du cours}

Le projet est explicitement conçu comme une synthèse, et il me semble important de
rendre cette intention visible dès l'introduction. Le tableau~\ref{tab:reinvest}
met en regard chaque chapitre du cours et la manière précise dont il se retrouve
mobilisé dans \ilang. Cette mise en correspondance n'est pas un simple exercice de
style. Elle constitue un fil conducteur que je m'efforce de suivre tout au long du
rapport, en signalant, à chaque fois que c'est pertinent, à quelle notion du cours
se rattache tel choix d'implémentation.

\begin{table}[ht]
\centering
\caption{Réinvestissement des chapitres du cours dans le projet.}
\label{tab:reinvest}
\small
\begin{tabular}{@{}p{4.8cm}p{9cm}@{}}
\toprule
\textbf{Chapitre du cours} & \textbf{Concept réinvesti dans \ilang} \\
\midrule
2, Assembleur & Pile d'appels et cadre \code{\%ebp}, reproduits pour les fonctions,
et générés tels quels par le back-end assembleur. \\
3, Programme C compilé & Niveaux d'optimisation et pliage de constantes. \\
4, Analyse lexicale et syntaxique & Analyseur lexical écrit avec Flex et parseur
écrit avec Bison. \\
5, Yacc/Bison & Grammaire LALR(1) et résolution, ou plutôt constat d'absence, des
conflits shift/reduce. \\
7, Génération de code & Arbre syntaxique abstrait, table des symboles, évaluation
récursive. \\
8, Optimisations & Pliage de constantes et élimination du code mort sur l'arbre. \\
9, Automates et regex & Tokens reconnus par l'automate fini engendré par Flex. \\
\bottomrule
\end{tabular}
\end{table}

\section{Structure du document}

Le rapport suit une progression que je veux la plus naturelle possible, allant du
général au particulier puis à la prise de recul. Le chapitre~\ref{chap:etat} dresse
un état de l'art des techniques d'interprétation et de compilation et explique
comment mon travail se positionne par rapport à elles. Le
chapitre~\ref{chap:conception} présente l'analyse des besoins, l'architecture que
j'ai retenue et la conception détaillée, en particulier la grammaire, l'arbre
syntaxique et les algorithmes principaux. Le chapitre~\ref{chap:implementation}
décrit l'implémentation proprement dite, module par module, avec de nombreux
extraits de code commentés et des exemples concrets. Le
chapitre~\ref{chap:validation} expose ma stratégie de test et les résultats
obtenus, y compris les mesures de performance et une étude de cas de bout en bout.
Le chapitre~\ref{chap:discussion} prend du recul, analyse de façon critique les
forces et les limites du projet, et revient en détail sur les difficultés que j'ai
rencontrées. Le chapitre~\ref{chap:conclusion} conclut et ouvre des perspectives.
Les annexes regroupent enfin l'intégralité du code source, les manuels
d'installation et d'utilisation, l'assembleur réellement généré et les données
expérimentales.
